\documentclass[../../../include/open-logic-section]{subfiles}

\begin{document}

\olfileid{sth}{z}{sep}
\olsection{分离}

我们首先引入一个替代朴素概括原则的原理：

\begin{axiom}[分离公理模式]
对于每个公式 $\phi(x)$，下述是一条公理：对任意 $A$，集合 $\Setabs{x \in A}{\phi(x)}$ 存在。
\end{axiom}

请注意，这不是单一的公理。它是一个公理\emph{模式}。存在\emph{无穷多条}分离公理；每个公式 $\phi(x)$ 对应一条。该模式同样可以（且通常）写为如下形式：

\begin{defish}
对于任何不包含“$S$”的公式 $\phi(x)$，下述是一条公理：
\[
	\forall A \exists S \forall x(x \in S \liff (\phi(x) \land x \in A)).
\]
\end{defish}

依照 \olref[sth][][]{part} 开头指出的约定，分离公理中的公式~$\phi$ 可以含有参数。\footnote{关于此含义的解释，请参见 \olref[sfr][infinite][induction]{natinductionschema} 之后的讨论。}

我们一直刻画的集合累积迭代观直接为分离公理提供了辩护。为了理解这一点，设 $A$ 是一个集合。那么 $A$ 在某个阶段~$S$ 形成（根据 \stageshier）。既然 $A$ 在阶段~$S$ 形成，则 $A$ 的所有成员都在阶段 $S$ 之前形成（根据 \stagesacc）。特别地，考虑所有既是 $A$ 的成员又满足 $\phi$ 的集合；显然，所有这些集合也都在阶段~$S$ 之前形成。因此，它们在阶段~$S$ 也能形成一个集合 $\Setabs{x \in A}{\phi(x)}$（根据 \stagesacc）。

与朴素概括不同，分离公理避免了 Russell悖论。因为我们不能简单地断言集合 $\Setabs{x}{x \notin x}$ 的存在性。相反，\emph{给定}某个集合~$A$，我们可以断言集合 $R_A = \Setabs{x \in A}{x \notin x}$ 的存在性。但这仅仅证明了 $R_A \notin R_A$ 且 $R_A \notin A$，这并没有什么令人担忧的。

然而，分离公理有一个直接而引人注目的推论：

\begin{thm}\ollabel{thm:NoUniversalSet}
不存在\emph{全集}，即 $\Setabs{x}{x = x}$ 不存在。
\end{thm}

\begin{proof}
归谬法，假设 $V$ 是一个全集。那么根据分离公理，$R =
\Setabs{x \in V}{x \notin x} = \Setabs{x}{x \notin x}$ 存在，这与Russell悖论矛盾。
\end{proof}

全集的不存在——或者确切地说，集合层次的开放性——是累积迭代观背后最基本的思想之一。因此值得留意，直观上我们可以通过另一条途径得出这一结论。一个全集必定是它自身的元素。但是，根据我们的累积迭代观，每个集合都是在它的所有元素形成之后紧接着的第一个阶段，（首次）出现在层次中的。但这意味着\emph{没有任何}集合是自属的。因为任何自属的集合，其首次出现都必须紧接在其自身首次出现的阶段之后，这显然是荒谬的。（我们将在 \olref[sth][spine][rank]{defnsetrank} 中看到，如何使用集合的“秩”概念使这一解释更加严谨。不过，为此我们需要先确立更多的公理。）

以下是分离公理和外延公理的更多推论。

\begin{prop}\ollabel{prop:emptyexists}
如果存在任何集合，则 $\emptyset$ 存在。
\end{prop}

\begin{proof}
如果 $A$ 是一个集合，那么根据分离公理，$\emptyset = \Setabs{x \in A}{x \neq x}$ 存在。
\end{proof}

\begin{prop}
对于任意集合 $A$ 和 $B$，$A \setminus B$ 存在。
\end{prop}

\begin{proof}
根据分离公理，$A \setminus B = \Setabs{x \in A}{x \notin B}$ 存在。
\end{proof}

此外，（几乎）任意的交集也存在：

\begin{prop}\ollabel{prop:intersectionsexist}
如果 $A \neq \emptyset$，则 $\bigcap A = \Setabs{x}{(\forall y \in A)x \in y}$ 存在。
\end{prop}

\begin{proof}
设 $A \neq \emptyset$，则存在某个 $c \in A$。那么 $\bigcap A = \Setabs{x}{(\forall y \in A)x \in y} = \Setabs{x \in c}{(\forall y \in A)x \in y}$，后者根据分离公理存在。
\end{proof}

但请注意条件 $A \neq \emptyset$；因为依空真，$\bigcap \emptyset$ 将成为全集，这与 \olref{thm:NoUniversalSet} 矛盾。

\end{document}